\section{Related Work}
\label{sec:related-work}

\paragraph{Automated LLM red teaming.}
\citet{perez2022redteaming} showed that language models can themselves generate
large and diverse red-team test sets, substantially increasing the scale of
harm discovery compared with purely manual adversarial authoring. Subsequent
systems work has pushed in a similar direction: \citet{radharapu2023aart}
propose AI-assisted data-generation recipes for building adversarial evaluation
sets for new LLM-powered applications, while \citet{mazeika2024harmbench}
argue that automated red teaming needs a standardized framework for comparing
attack methods and robust refusals. \system\ sits closer to this systems and
infrastructure tradition than to single-attack papers: its contribution is not
a new jailbreak search algorithm, but an execution and reporting stack that can
run campaigns, retain traces, evaluate them, and surface actionable outputs.

\paragraph{Jailbreaks and refusal bypass.}
The contemporary jailbreak literature documents a wide spectrum of alignment
failures. \citet{wei2023jailbroken} frame jailbreak success in terms of safety
training failure modes, while \citet{zou2023universal} show that automatically
optimized adversarial suffixes can transfer across aligned models. Empirical
systematization work also finds that semantically meaningful jailbreak prompts
remain widely effective and can be assisted by automation
\citep{yu2024dontlisten}. Our archived Mistral traces belong to this broader
family of failures: they are not optimized suffix attacks, but they do show
that plain-English override and persona-hijack prompts can still produce clear
alignment breakdowns.

\paragraph{Prompt injection and LLM-integrated systems.}
For LLM applications embedded in larger software stacks, prompt injection is
particularly important. \citet{greshake2023notwhat} demonstrate indirect prompt
injection against real-world LLM-integrated applications, while
\citet{toyer2024tensortrust} construct a large prompt-injection dataset and
benchmark from the Tensor Trust game. Those papers motivate \system's attention
to tool misuse, system prompt extraction, context poisoning, and other
application-facing failure modes rather than harmful-completion prompts alone.

\paragraph{Evaluator-grounded safety scoring.}
As attack papers proliferated, so did exaggerated claims of success.
\citet{souly2024strongreject} argue that many jailbreak evaluations overestimate
severity and propose \sr\ as a stricter benchmark and scoring rubric based on
refusal, convincingness, and specificity. \citet{mazeika2024harmbench} make a
parallel case for standardized evaluation in automated red teaming. In this
paper we use \sr\ as an auxiliary scoring layer over the repository's existing
ATT\&CK-inspired campaign probes, not as a substitute benchmark. That design is
useful for comparing responses within this framework, but it should not be read
as a direct leaderboard comparison to HarmBench or public \sr\ benchmark
results.

\paragraph{Systematic and socio-technical perspectives.}
Recent 2026 work broadens red teaming beyond attack-success reporting alone.
\citet{srivastava2026systematicredteaming} provide a broad systematic review of
algorithmic red-teaming methods for scalable assurance workflows, while
\citet{seghid2026emergingthreats} situate red teaming inside a wider misuse and
governance taxonomy spanning security, disinformation, and human-harm domains.
Human-centered analysis also matters for methodological validity:
\citet{alvaradogarcia2026sociotechnical} show that practitioner assumptions and
dataset-construction practices can introduce blind spots in what risks are
measured. Complementary accessibility-focused TEVV work
\citep{waters2026accessibilitytevv} argues that adversarial evaluation should
explicitly include disability- and access-related failure modes rather than
treating them as out-of-scope edge cases.

\paragraph{Scaling and FAR.AI-adjacent empirical work.}
Recent work has also asked whether larger models become more, rather than less,
susceptible to certain safety failures. \citet{bowen2024scalingtrends} examine
data-poisoning vulnerability across model sizes and report stronger
susceptibility at larger scales in that setting. Our repository includes model
registry support and plotting for analogous comparisons, but the executed sample
available here is too narrow to sustain a broad scaling claim. We therefore use
their work as conceptual motivation, while treating our own size figure as
exploratory.

\paragraph{Recent FAR.AI model-evaluation line.}
FAR.AI's recent sequence of evaluations has covered complementary parts of the
red-teaming surface: data-poisoning and jailbreak-tuning susceptibility
\citep{bowen2024scalingtrends,murphy2025jailbreaktuning}; hidden capability
exposure via the safety-gap lens \citep{dombrowski2025safetygap}; benchmark
pressure from more difficult harmful-query datasets
\citep{hollinsworth2025clearharm}; model-to-model movement on persuasive-harm
behavior \citep{kowal2026persuasion}; and newly systematic analysis of
prefill-level jailbreak vectors in open-weight systems
\citep{struppek2026prefill}. Taken together, this body of work motivates
continuous, technique-adaptive red teaming rather than one-off evaluations tied
to a single jailbreak style or benchmark snapshot.

\paragraph{Ecosystem trends and evaluation urgency.}
Over the last year, empirical organizations including FAR.AI, METR, and Epoch
AI have expanded public evidence on frontier-model capability growth, evaluation
practice, and deployment-relevant risk measurement
\citep{bowen2024scalingtrends,metr2025deepseekqwen,epoch2025impact}. In
parallel, frontier deployments are increasingly coupled with staged access and
capability-gated release workflows, which raises the bar for red-teaming
operations \citep{openai2026trustedcyber,anthropic2026glasswing}. At the same
time, providers continue to revise and streamline safety-governance mechanisms;
public commentary sometimes interprets these shifts as reducing robustness,
which creates additional uncertainty for evaluators and downstream governance
loops. Additionally, governance remains fragmented across jurisdictions and
firms, even as a small number of frontier labs retain disproportionate
influence over model release norms, access controls, and de facto safety
baselines. In the EU policy cycle around the AI Act, contemporaneous reporting
described Mistral as actively pushing to narrow foundation-model obligations,
illustrating how concentrated industry power can shape the practical governance
environment for red teaming
\citep{lomas2023mistraleuai,goldman2023euaiactmistral}. This reinforces a practical conclusion: red teaming should be
continuous, technique-adaptive, and broad in scope, including domains often
treated as lower-priority (for example misinformation/disinformation pathways,
mental-health-adjacent interactions, and sycophancy-like behavioral failures),
because those behaviors can still surface policy-relevant weaknesses and shape
future evaluation frameworks. Related agentic-system analyses
\citep{radanliev2026agenticthreats} further support expanding threat models
beyond single-turn prompt perturbations toward executive-layer and closed-loop
failure modes.
